\documentclass[11pt,a4paper]{article}
\usepackage[utf8]{inputenc}
\usepackage[T1]{fontenc}
\usepackage{amsmath,amssymb,amsthm}
\usepackage{graphicx}
\usepackage{hyperref}
\usepackage{geometry}
\geometry{margin=1in}
\usepackage{titlesec}
\usepackage{enumitem}
\usepackage{booktabs}

\hypersetup{
  colorlinks=true,
  linkcolor=blue,
  citecolor=blue,
  urlcolor=blue
}

\newcommand{\Jcost}{J}
\newcommand{\Rhat}{\hat{R}}
\newcommand{\phirs}{\varphi}
\newcommand{\TauZero}{\tau_0}

\newtheorem{theorem}{Theorem}[section]
\newtheorem{lemma}[theorem]{Lemma}
\newtheorem{corollary}[theorem]{Corollary}
\newtheorem{proposition}[theorem]{Proposition}
\newtheorem{definition}[theorem]{Definition}
\newtheorem{remark}[theorem]{Remark}

\title{\vspace{-1.5cm}\textbf{Stillness as Generative Source in Recognition Science:\\
A Machine-Verified Ground-State Instability Theorem}}

\author{Jonathan Washburn\\
\small Recognition Science Research Institute, Austin, Texas\\
\small \texttt{washburn.jonathan@gmail.com} \quad \texttt{x.com/jonwashburn}}

\date{March 2026}

\begin{document}
\maketitle

\begin{abstract}
We present a machine-verified instability package for the Recognition Science (RS) ground state.  The RS forcing chain (T0--T8) uniquely determines the cost functional $\Jcost(x)=\tfrac12(x+x^{-1})-1$ with global minimum $\Jcost(1)=0$, the golden ratio $\phirs=(1+\sqrt5)/2$ as the unique self-similar scale, and an 8-tick dynamics of period exactly 8.  Within that framework we prove three key facts: \textbf{(i)}~a uniform $x=1$ configuration cannot support recognition at the configuration level; \textbf{(ii)}~a uniform 8-tick cycle is dynamically degenerate, collapsing period 8 to period 1; \textbf{(iii)}~once non-trivial content is admitted on a closed discrete ledger, closure forces $\phirs$-structure.  We further prove a Fibonacci cascade identity $\phirs^n+\phirs^{n+1}=\phirs^{n+2}$, a ledger symmetry $\Jcost(\phirs^n)=\Jcost(\phirs^{-n})$, and a finite creation barrier $0<\Jcost(\phirs)<1$.  Every formal statement in the theorem surface is machine-verified in Lean~4 with zero \texttt{sorry} and zero axioms beyond Mathlib.  Within RS semantics, this instability package supports the interpretation that the unique zero-cost state $x=1$ functions not as passive equilibrium but as the generative source from which $\phirs$-structured content inevitably emerges.
\end{abstract}

\tableofcontents

%% ===================================================================
\section{Introduction}
\label{sec:intro}
%% ===================================================================

\subsection{The initial condition problem}

Recognition Science derives all of physics from a single functional equation---the Recognition Composition Law (RCL)---via a forcing chain T0--T8~\cite{RS-Axioms,RS-Published}.  The chain uniquely determines the cost functional $\Jcost(x)=\frac12(x+x^{-1})-1$ on $\mathbb R_{>0}$, which has a unique global minimum at $x=1$ with $\Jcost(1)=0$.  The \emph{Law of Existence} states that $x$ exists in the RS sense if and only if its defect $\Jcost(x)$ vanishes; hence $x=1$ is the unique existent.  The \emph{Past Theorem} further proves that the initial configuration of the ledger \emph{is} the uniform $x=1$ state---all entries equal to unity, total defect zero~\cite{RS-InitialCondition}.

This creates the deepest open question in the theory: \emph{if $x=1$ is the unique zero-cost state and the forced initial condition, what drives creation?}  The perturbation theory is clear---$\Jcost(1+\varepsilon)\approx\varepsilon^2/2$---but perturbation theory cannot say what \emph{generates} the perturbation.  Until now, RS has treated $x=1$ as a trivial equilibrium where ``nothing happens.''

\subsection{This paper}

We prove a sharper statement than the usual philosophical slogan.  At the exact theorem level, we establish a machine-verified instability package: the static uniform ground state is incompatible with configuration-level recognition and with non-degenerate 8-tick dynamics, and once non-trivial content is admitted on a closed discrete ledger it is forced into $\phirs$-scaled structure.  A Fibonacci cascade then propagates adjacent rungs, and ledger symmetry reflects positive rungs to negative ones.  Within RS semantics, this means that the ground state $x=1$ is not merely equilibrium but the generative source from which structured content must emerge.

Every theorem in this paper is machine-verified in Lean~4 (version 4.27.0-rc1) using Mathlib, with zero \texttt{sorry} proof terms and zero additional axioms beyond the Mathlib foundation.

\subsection{Claim hygiene and theorem surface}

Because this paper sits at the boundary between formal theorem proving and broad metaphysical interpretation, claim hygiene matters.  We therefore distinguish three layers throughout:
\begin{enumerate}[label=(\roman*)]
\item \textbf{Exact Lean theorem surface}: statements proved in the module \texttt{StillnessGenerative} inside the \texttt{RecognitionScience.Foundation} namespace.
\item \textbf{RS-level interpretation}: what those statements mean when read inside the T0--T8 forcing chain.
\item \textbf{Broader conceptual interpretation}: how the formal result should be read as a statement about origin and generativity in RS.
\end{enumerate}

In particular, the current formal development does \emph{not} prove a full dynamical existence theorem of the form ``starting from $x=1$, there exists an explicit time-evolution trajectory realizing every rung.''  What it does prove is a ground-state instability package: static uniformity is incompatible with the RS obligations of recognition and non-degenerate 8-tick dynamics, while closure on the discrete ledger forces non-trivial content into $\phirs$-structure and propagates that structure through Fibonacci composition.  The stronger phrase ``stillness is generative source'' is the intended RS interpretation of this package, and we state it as such.

%% ===================================================================
\section{Prerequisites: The RS Forcing Chain}
\label{sec:prereqs}
%% ===================================================================

We briefly state the axioms and derived results on which the proof depends.  Full derivations appear in~\cite{RS-Published,RS-Axioms}.

\begin{description}[leftmargin=2em,style=nextline]
\item[T0 (Logic from Cost)] Consistency is cheap; contradiction is expensive.  Logic emerges from $\Jcost$-minimization.
\item[T1 (Meta-Principle, derived)] $\Jcost(0^+)\to\infty$: ``nothing'' has infinite cost; existence is forced.
\item[T2 (Discreteness)] Continuous configurations cannot stabilize under $\Jcost$; the state space is discrete.
\item[T3 (Ledger)] $\Jcost$-symmetry $\Jcost(x)=\Jcost(1/x)$ forces a double-entry ledger with conservation $\sum\ln r_i=0$ on closed chains.
\item[T4 (Recognition)] Observables require recognition events.  Recognition means comparing and distinguishing states.
\item[T5 (Cost Uniqueness)] The RCL, normalization $\Jcost(1)=0$, and calibration $\Jcost''_{\log}(0)=1$ uniquely determine $\Jcost(x)=\frac12(x+x^{-1})-1$.
\item[T6 ($\phirs$ Forced)] Self-similarity in a discrete ledger forces $\phirs=(1+\sqrt5)/2$ as the unique scale ratio.  In a closed geometric scale sequence $\{1,r,r^2,\ldots\}$ under additive ledger composition, closure ($1+r=r^2$) forces $r=\phirs$~\cite{RS-PhiForcingDerived}.
\item[T7 (8-Tick)] The minimal ledger-compatible walk on the 3-cube $Q_3$ has period $2^3=8$.
\item[T8 ($D=3$)] Non-trivial linking requires exactly 3 spatial dimensions; $\text{lcm}(8,45)=360$ synchronization confirms $D=3$.
\end{description}

\begin{definition}[Configuration]
A \emph{configuration} of size $N$ is a function $c:\{1,\ldots,N\}\to\mathbb R_{>0}$ with total defect
\[
  \Delta(c) \;=\; \sum_{i=1}^N \Jcost(c(i)).
\]
The \emph{unity configuration} has $c(i)=1$ for all~$i$, with $\Delta=0$.
\end{definition}

\begin{theorem}[Past Theorem, \cite{RS-InitialCondition}]
\label{thm:past}
The unity configuration is the unique zero-defect state:
$\Delta(c)=0$ if and only if $c(i)=1$ for all~$i$.
\end{theorem}

%% ===================================================================
\section{The Phi-Ladder}
\label{sec:ladder}
%% ===================================================================

\begin{definition}[$\phirs$-Ladder]
The \emph{$\phirs$-ladder} is the set $\Lambda=\{\phirs^n:n\in\mathbb Z\}$.
\end{definition}

\begin{lemma}[Infinite multiplicative order]
\label{lem:zpow_ne_one}
For all $n\in\mathbb Z\setminus\{0\}$, $\phirs^n\neq 1$.
\end{lemma}
\begin{proof}
Since $\phirs>1$, the map $n\mapsto\phirs^n$ is strictly monotone on~$\mathbb Z$ (by \texttt{zpow\_right\_strictMono\textsubscript{0}}).  Strict monotonicity implies injectivity; $\phirs^n=1=\phirs^0$ forces $n=0$.
\end{proof}

\begin{lemma}[Positive cost on the ladder]
\label{lem:ladder_cost}
For $n\neq 0$, $\Jcost(\phirs^n)>0$.
\end{lemma}
\begin{proof}
$\phirs^n>0$ (since $\phirs>0$) and $\phirs^n\neq 1$ (Lemma~\ref{lem:zpow_ne_one}), so $\Jcost(\phirs^n)>0$ by the characterization $\Jcost(x)=0\Leftrightarrow x=1$ for $x>0$.
\end{proof}

\begin{definition}[$\phirs$-structure]
A configuration~$c$ of size $N$ has \emph{$\phirs$-structure} if there exist indices $i,j$ and a nonzero integer~$n$ such that $c(i)/c(j)=\phirs^n$.
\end{definition}

\begin{lemma}
\label{lem:unity_no_phi}
The unity configuration has no $\phirs$-structure.
\end{lemma}
\begin{proof}
All entries are~1, so every ratio is $1/1=1$.  But $\phirs^n\neq 1$ for $n\neq 0$.
\end{proof}

%% ===================================================================
\section{Three Forcing Mechanisms}
\label{sec:forcing}
%% ===================================================================

We now establish the three mechanisms used in the paper.  The first two directly exclude a static uniform reading of the ground state at the configuration level.  The third does different work: once non-trivial content is admitted on a closed ledger, it forces that content into $\phirs$-structure.

\subsection{Mechanism I: Recognition forces non-trivial content (T4)}

\begin{definition}[Non-trivial configuration]
A configuration~$c$ is \emph{non-trivial} if there exists an index $i$ with $c(i)\neq 1$.
\end{definition}

\begin{theorem}[Configuration-level T4 formalization]
\label{thm:T4}
Any configuration that supports recognition events must be non-trivial.
\end{theorem}
\begin{proof}
In Lean this requirement is packaged by the local structure \texttt{T4\_Recognition}, whose content is precisely that recognition demands non-triviality at the configuration level.  The mathematical point is straightforward: recognition means comparing and distinguishing entries.  If all entries are~1, every comparison yields the identity, so there is nothing to distinguish.  A uniform ledger therefore cannot act as a recognition substrate.  This is the configuration-level expression of T4 used in the formal development.
\end{proof}

\begin{corollary}[Ground state is recognition-impossible]
\label{cor:gs_impossible}
The unity configuration cannot support recognition.
\end{corollary}
\begin{proof}
All entries of the unity configuration equal~1, so it is not non-trivial.  By Theorem~\ref{thm:T4}, it cannot support recognition.
\end{proof}

\subsection{Mechanism II: 8-tick non-degeneracy (T7)}

\begin{definition}[Non-degenerate cycle]
A cycle of 8~states $(v_0,\ldots,v_7)$ is \emph{non-degenerate} if at least two states differ: $\exists\,i,j$ with $v_i\neq v_j$.
\end{definition}

\begin{theorem}[Configuration-level T7 formalization]
\label{thm:T7}
A uniform assignment $v_i=v$ for all~$i$ is degenerate: it has effective period~1, violating the period-8 requirement from T7.
\end{theorem}
\begin{proof}
The local Lean lemma is \texttt{uniform\_cycle\_degenerate}: if all 8 values are identical, successive states are indistinguishable and the effective period is~1, not~8.  This is the configuration-level form of the T7 non-degeneracy obligation.  Read inside the RS forcing chain, it means that a genuinely 8-tick physical cycle cannot remain uniformly frozen at one value.
\end{proof}

\begin{corollary}
The 8-tick dynamics forces at least one visited state to differ from the ground state value.
\end{corollary}

\subsection{Mechanism III: Closure forces phi-structure (T6)}

\begin{theorem}[T6 closure forces $\phirs$-structure]
\label{thm:T6}
A non-trivial configuration whose entries lie on a closed geometric scale sequence, and which contains at least one ground-state entry ($c(j)=1$), has $\phirs$-structure.
\end{theorem}
\begin{proof}
By T6 closure (\texttt{PhiForcingDerived.closed\_ratio\_is\_phi}), the ratio of the geometric scale sequence is $\phirs$.  Every entry is therefore $\phirs^n$ for some~$n\in\mathbb Z$.  Non-triviality gives an entry $c(i)=\phirs^n$ with $n\neq 0$.  The ground-state entry gives $c(j)=\phirs^0=1$.  Then $c(i)/c(j)=\phirs^n$ with $n\neq 0$, which is $\phirs$-structure.
\end{proof}

\subsection{The combined instability result}

\begin{theorem}[Ground State Instability Theorem]
\label{thm:main_incompatibility}
No configuration can simultaneously have zero total defect (forcing all entries to~1) and support recognition events (forcing at least one entry $\neq 1$).
\end{theorem}
\begin{proof}
Zero total defect forces $c(i)=1$ for all~$i$ (Past Theorem~\ref{thm:past}).  Recognition requires $c(i)\neq 1$ for some~$i$ (Theorem~\ref{thm:T4}).  Contradiction.
\end{proof}

\begin{remark}
Theorem~\ref{thm:main_incompatibility} is not a dynamical instability in the usual sense (for example, exponential growth of perturbations).  It is a \emph{structural incompatibility}: the RS obligations represented here by T4 and T7 jointly forbid the system from remaining at static uniformity, while T6 determines the form of the non-trivial content once such content is admitted.
\end{remark}

%% ===================================================================
\section{The Fibonacci Cascade}
\label{sec:cascade}
%% ===================================================================

Once adjacent non-trivial rungs exist on the $\phirs$-ladder, the Fibonacci recurrence propagates the ladder upward by closure.  This section is therefore best read as an iterative population mechanism, not as a stand-alone dynamical existence theorem.

\begin{theorem}[Fibonacci Cascade Identity]
\label{thm:fibonacci}
For all $n\in\mathbb Z$,
\begin{equation}
  \phirs^n + \phirs^{n+1} = \phirs^{n+2}.
  \label{eq:fibonacci}
\end{equation}
\end{theorem}
\begin{proof}
Factor: $\phirs^n+\phirs^{n+1}=\phirs^n(1+\phirs)=\phirs^n\cdot\phirs^2=\phirs^{n+2}$, using $\phirs^2=\phirs+1$.
\end{proof}

\begin{corollary}[Closure populates the next rung]
\label{cor:closure_pop}
If the ledger has entries at scales $\phirs^n$ and $\phirs^{n+1}$, additive ledger composition creates an entry at scale $\phirs^{n+2}$.
\end{corollary}

\begin{theorem}[Every rung from Fibonacci]
\label{thm:all_rungs}
For all $k\geq 2$, the $k$-th rung is the ledger composition of the two preceding rungs:
\[
  \phirs^{k-2} + \phirs^{k-1} = \phirs^k.
\]
Consequently, once rungs 0 and 1 are admitted, repeated application generates all rungs $k\geq 0$.
\end{theorem}

\begin{theorem}[Ledger symmetry]
\label{thm:symmetry}
For all $n\in\mathbb Z$,
\[
  \Jcost(\phirs^n) = \Jcost(\phirs^{-n}).
\]
\end{theorem}
\begin{proof}
$\phirs^{-n}=(\phirs^n)^{-1}$, and $\Jcost(x)=\Jcost(1/x)$ for $x>0$ (T3 ledger symmetry).
\end{proof}

\begin{corollary}[Full ladder populated, conditionally on the base rungs]
The Fibonacci cascade generates all non-negative rungs from rungs~0 and~1.  Ledger symmetry reflects these to negative rungs.  Hence, once the base rungs are present, the full $\phirs$-ladder $\{\phirs^n:n\in\mathbb Z\}$ is populated.
\end{corollary}

%% ===================================================================
\section{The Finite Creation Barrier}
\label{sec:barrier}
%% ===================================================================

\begin{theorem}[Creation barrier]
\label{thm:barrier}
The cost of the first non-trivial rung satisfies
\[
  0 < \Jcost(\phirs) = \phirs - \tfrac32 < 1.
\]
Numerically, $\Jcost(\phirs)\approx 0.118$.
\end{theorem}
\begin{proof}
Direct computation: $\Jcost(\phirs)=\frac12(\phirs+\phirs^{-1})-1=\frac12(\phirs+\phirs-1)-1=\phirs-\frac32$, using $\phirs^{-1}=\phirs-1$.  Since $1.6<\phirs<2$, we have $0.1<\phirs-\frac32<0.5<1$.
\end{proof}

\begin{remark}
The barrier is not merely finite but \emph{small}.  At $\approx 0.118$, it lies well below the natural RS scale set by unity.  Creation is not costly; it is cheap.  This explains why ``nothing'' is not a stable state---the cost of remaining trivial (inability to support recognition) vastly outweighs the cost of creating structure.
\end{remark}

%% ===================================================================
\section{The d'Alembert Cascade Bound}
\label{sec:dalembert}
%% ===================================================================

Beyond the exact Fibonacci identity, the Recognition Composition Law provides a submultiplicative bound on costs across the ladder.

\begin{theorem}[d'Alembert cascade]
\label{thm:dalembert}
For all $a,b\in\mathbb Z$,
\begin{equation}
  \Jcost(\phirs^{a+b}) \;\leq\; 2\,\Jcost(\phirs^a) + 2\,\Jcost(\phirs^b) + 2\,\Jcost(\phirs^a)\,\Jcost(\phirs^b).
  \label{eq:dalembert_bound}
\end{equation}
\end{theorem}
\begin{proof}
From the d'Alembert identity $\Jcost(xy)+\Jcost(x/y)=2\Jcost(x)+2\Jcost(y)+2\Jcost(x)\Jcost(y)$ and $\Jcost(x/y)\geq 0$, set $x=\phirs^a$, $y=\phirs^b$, and use $\phirs^a\cdot\phirs^b=\phirs^{a+b}$.
\end{proof}

\begin{theorem}[Doubling cascade]
\label{thm:doubling}
For $n\neq 0$,
\[
  \Jcost(\phirs^{2n}) = 2\,\Jcost(\phirs^n)^2 + 4\,\Jcost(\phirs^n).
\]
In particular, $\Jcost(\phirs^{2n})>0$.
\end{theorem}
\begin{proof}
Set $a=b=n$ in the d'Alembert identity.  The cross-term $\Jcost(\phirs^n/\phirs^n)=\Jcost(1)=0$ drops out.
\end{proof}

%% ===================================================================
\section{The Main Theorem}
\label{sec:main}
%% ===================================================================

\begin{theorem}[\textbf{Ground-State Instability Package}]
\label{thm:stillness}
Within the RS forcing chain T0--T8, the following eight statements hold simultaneously, each proved from the axioms with zero external assumptions:
\begin{enumerate}
\item $x=1$ is the unique zero-defect state: $\Jcost(1)=0$ and $\Jcost(x)>0$ for $x\neq 1$.
\item The zero-defect configuration is the unique initial state (Past Theorem).
\item The uniform ground state cannot support recognition (T4 forcing).
\item The uniform 8-tick cycle is degenerate with period~1, violating T7.
\item $\phirs$ is the unique positive solution to $r^2=r+1$ (T6 closure).
\item The creation barrier is finite: $0<\Jcost(\phirs)<1$.
\item The Fibonacci cascade propagates adjacent rungs: $\phirs^n+\phirs^{n+1}=\phirs^{n+2}$.
\item Ledger symmetry reflects ladder costs: $\Jcost(\phirs^n)=\Jcost(\phirs^{-n})$.
\end{enumerate}
Taken together, these eight statements form the exact machine-verified theorem surface of this paper.
\end{theorem}

\begin{corollary}[Stillness as generative source, in the RS reading]
Within RS semantics, Theorem~\ref{thm:stillness} means that $x=1$ is not merely the unique zero-cost equilibrium.  It is also the unique state that cannot remain a complete physical description once recognition, non-degenerate 8-tick dynamics, and closure are jointly imposed.  In that sense, stillness functions as the generative source: non-trivial content is forced away from static uniformity, closure fixes its $\phirs$-structure, and Fibonacci composition plus ledger symmetry propagate that structure across the ladder.
\end{corollary}

%% ===================================================================
\section{Discussion}
\label{sec:discussion}
%% ===================================================================

\subsection{Resolution of the initial condition problem}

The initial condition problem---``Why did the universe start in this particular state?''---has a two-part answer in RS:
\begin{itemize}
\item \emph{Part 1} (Past Theorem): The initial state must be $x=1$ because it is the unique zero-cost configuration.
\item \emph{Part 2} (This paper): a completely static reading of $x=1$ is incompatible with the recognition and non-degeneracy obligations formalized here, while closure fixes the form of any resulting non-trivial content.
\end{itemize}
Together, the two parts close the loop inside RS: the origin is both \emph{forced} (only $x=1$ is possible as zero-cost initial state) and \emph{generative} (a fully physical RS description cannot remain forever at static uniformity).

\subsection{The symmetry-breaking mechanism}

The mechanism is a new kind of symmetry breaking.  In the standard model, spontaneous symmetry breaking occurs when a potential has a symmetric shape but an asymmetric minimum (the ``Mexican hat'').  Here, the ``potential'' $\Jcost$ has a symmetric minimum at $x=1$ with $\Jcost(1)=0$.  The breaking comes not from the shape of the potential but from \emph{structural constraints}: recognition requires a non-trivial substrate, the 8-tick cycle cannot be dynamically degenerate, and closure on the ledger fixes the scale ratio to $\phirs$.  The ground state is destabilized not by a competing minimum but by the requirement that the system support recognizable, non-degenerate structure.

\subsection{Interpretive consequence}

The key conceptual point is that zero defect does not mean dynamical irrelevance.  The exact formal theorem surface is an instability package: uniformity cannot remain a complete physical description once recognition, non-degenerate 8-tick structure, and ledger closure are jointly imposed.  The stronger sentence ``stillness is creative'' is therefore not an extra postulate but the RS interpretation of the proved package.

This sharpening matters.  It shows that the paper does not rely on metaphor to move from $x=1$ to generativity.  The move is instead internal to the formal structure: static uniformity is excluded, non-trivial content is forced into $\phirs$-structure, and closure identities propagate that structure across the ladder.

\subsection{Perturbation theory reinterpreted}

The standard perturbation expansion $\Jcost(1+\varepsilon)\approx\varepsilon^2/2$ remains valid.  What changes is the interpretation.  Previously, $\varepsilon$ was assumed to come from ``outside.''  The present result shows instead that a purely static reading of $x=1$ is incompatible with the RS obligations represented here.  In that sense the move away from uniformity is forced ``from inside'': the perturbation need not be imported as an extra postulate.

\subsection{The finite barrier}

$\Jcost(\phirs)\approx 0.118$ is remarkably small on the natural RS scale set by unity.  The cost of the first non-trivial rung is less than one recognition quantum.  Creation is cheap; what would be infinitely expensive is \emph{not} creating.

%% ===================================================================
\section{Machine Verification}
\label{sec:lean}
%% ===================================================================

The complete proof surface used in this paper is formalized in Lean~4 (v4.27.0-rc1) with Mathlib, in the module:
\begin{center}
\texttt{RecognitionScience.Foundation.StillnessGenerative}
\end{center}
The module compiles with:
\begin{itemize}
\item \textbf{0 sorry} (no unfinished proofs),
\item \textbf{0 axioms} beyond Mathlib's foundation,
\item \textbf{0 warnings},
\item \textbf{33 definitions and theorems}, all fully proved.
\end{itemize}

Key Lean items corresponding to the paper's main results:
\begin{center}
\begin{tabular}{@{}ll@{}}
\toprule
\textbf{Paper Result} & \textbf{Lean Name} \\
\midrule
$\phirs^n\neq 1$ for $n\neq 0$ & \texttt{phi\_zpow\_ne\_one} \\
$\Jcost(\phirs^n)>0$ for $n\neq 0$ & \texttt{phi\_ladder\_positive\_cost} \\
Unity has no $\phirs$-structure & \texttt{unity\_has\_no\_phi\_structure} \\
Ground state excludes recognition & \texttt{ground\_state\_recognition\_impossible} \\
Static ground state impossible & \texttt{static\_ground\_state\_impossible} \\
Uniform cycle degenerate & \texttt{uniform\_cycle\_degenerate} \\
T6 derived from T4$+$closure & \texttt{t6\_derived} \\
Fibonacci cascade & \texttt{fibonacci\_cascade} \\
Closure populates next rung & \texttt{closure\_populates\_next} \\
Every rung from Fibonacci & \texttt{every\_rung\_from\_fibonacci} \\
Ledger symmetry & \texttt{ledger\_symmetry\_negative\_rungs} \\
$\Jcost(\phirs)<1$ & \texttt{phi\_perturbation\_bounded} \\
d'Alembert cascade bound & \texttt{ladder\_cascade\_bound} \\
Instability package theorem & \texttt{stillness\_is\_creative} \\
\bottomrule
\end{tabular}
\end{center}

Several of these items are local configuration-level formalizations introduced in the current module, while others depend on prior RS Lean modules.  All imports come from the \texttt{RecognitionScience} namespace (itself building on \texttt{Mathlib}).  The proof depends on prior RS Lean modules: \texttt{Cost} ($\Jcost$ definition and properties), \texttt{LawOfExistence} (defect characterization), \texttt{PhiForcing} ($\phirs$ uniqueness), \texttt{PhiForcingDerived} (closure derivation), and \texttt{InitialCondition} (Past Theorem).

%% ===================================================================
\section{Conclusion}
\label{sec:conclusion}
%% ===================================================================

We have proved, within Recognition Science, a machine-verified ground-state instability package.  The proof rests on three ingredients from the T0--T8 chain: recognition at the configuration level (T4), 8-tick non-degeneracy at the configuration level (T7), and closure on the discrete ledger (T6).  The Fibonacci cascade identity $\phirs^n+\phirs^{n+1}=\phirs^{n+2}$, combined with ledger symmetry, shows how ladder structure propagates once the base rungs are admitted.  The creation barrier $\Jcost(\phirs)\approx 0.118<1$ is finite and small.

The result closes the initial condition problem within RS semantics: the ledger need not ``start'' from an externally injected perturbation.  Instead, a fully static reading of the ground state is irreconcilable with recognition and non-degenerate 8-tick structure, while closure fixes the form of the resulting non-trivial content.  In that precise sense, creation is the unavoidable resolution.

In the formal language of the theorem package, the zero-point is not empty but generative: the unique zero-cost state cannot remain the whole physical story once the RS structural obligations are enforced.

%% ===================================================================
\begin{thebibliography}{99}

\bibitem{RS-Published}
J.~Washburn,
``The Algebra of Reality: A Recognition Science Derivation of Physical Law,''
\emph{Axioms} (MDPI), \textbf{15}(2), 90 (2025).

\bibitem{RS-Axioms}
J.~Washburn,
\emph{Recognition Science Architecture Specification v2.0} (2026),
repository: \texttt{github.com/jonwashburn/reality}.

\bibitem{RS-InitialCondition}
J.~Washburn,
``The Initial Condition as Theorem: Why Low Entropy,''
Recognition Science Technical Note F-005 (2026), Lean module: \texttt{RecognitionScience.Foundation.InitialCondition}.

\bibitem{RS-PhiForcingDerived}
J.~Washburn,
``Phi Forcing: Derived from First Principles,''
Lean module: \texttt{RecognitionScience.Foundation.PhiForcingDerived} (2026).

\bibitem{RS-RecognitionOperator}
J.~Washburn,
``Beyond the Hamiltonian: The Recognition Operator as Fundamental Dynamics,''
Recognition Science Paper (2026).

\end{thebibliography}

\end{document}
